% EXPERIMENTAL METHODOLOGY
% Last generated: 2026-01-15T05:34:34.872Z

\subsection{Experimental Design}

We employ Monte Carlo simulation with systematic parameter sweeps to investigate our research questions. Each experiment configuration is run multiple times with different random seeds to establish statistical distributions.

\subsubsection{Parameter Sweeps}

For each research question, we vary a single parameter while holding others constant (ceteris paribus):

\begin{table}[htbp]
\centering
\caption{Experimental parameter configurations}
\label{tab:experiments}
\small
\begin{tabular}{@{}llll@{}}
\toprule
Experiment & Swept Parameter & Range & Runs \\
\midrule
00 Baseline & (none) & -- & 2{,}000 \\
01 Calibration & participation targets (grid) & $3\times3$ & 270 \\
02 Ablation & governance mechanisms (grid) & $2\times2\times2$ & 200 \\
03 Quorum Sensitivity & \texttt{quorumPercentage} & 0.01--0.10 & 150 \\
04 Governance Capture & mitigation configs (grid) & 27 & 675 \\
05 Proposal Pipeline & pipeline params (grid) & 9 & 270 \\
06 Treasury Resilience & treasury params (grid) & 12 & 300 \\
07 Inter-DAO & scenarios & 5 & 300 \\
08 Scale Sweep & \texttt{totalMembers} & 50--500 & 150 \\
09 Voting Mechanisms & mechanism $\times$ quorum (grid) & 6 & 180 \\
\bottomrule
\end{tabular}
\end{table}

\begin{table}[htbp]
\centering
\caption{Mapping of research questions to experiments}
\label{tab:rq-mapping}
\begin{tabular}{@{}ll@{}}
\toprule
Research Question & Experiments \\
\midrule
RQ1: Participation dynamics & 01 Calibration, 03 Quorum Sensitivity \\
RQ2: Governance capture & 04 Governance Capture \\
RQ3: Proposal pipeline & 05 Proposal Pipeline \\
RQ4: Treasury resilience & 06 Treasury Resilience \\
RQ5: Inter-DAO cooperation & 07 Inter-DAO \\
Supporting & 00 Baseline, 02 Ablation, 08 Scale, 09 Voting \\
\bottomrule
\end{tabular}
\end{table}

\subsubsection{Baseline Configuration}

All experiments use a common baseline derived from Compound governance:

\begin{itemize}
    \item Members: 200 (unless swept)
    \item Token distribution: Power-law ($\alpha = 1.5$)
    \item Quorum: 4\% (unless swept)
    \item Voting period: 7 days
    \item Timelock: 2 days
    \item Proposal frequency: 0.5/day
    \item Simulation length: 2{,}000 steps (days)
\end{itemize}

\subsection{Metrics}

We capture the following metrics for each simulation run:

\subsubsection{Primary Metrics}

\begin{description}
    \item[Proposal Pass Rate] Fraction of proposals that achieved quorum and majority support
    \item[Average Turnout] Mean participation rate across all votes
    \item[Final Gini] Gini coefficient of token holdings at simulation end
    \item[Treasury Efficiency] Funds deployed to successful projects vs. total treasury
\end{description}

\subsubsection{Secondary Metrics}

\begin{description}
    \item[Quorum Failure Rate] Proposals failing specifically due to quorum
    \item[Delegation Concentration] Herfindahl-Hirschman Index of delegated power
    \item[Time to Decision] Average steps from proposal to resolution
\end{description}

\subsection{Statistical Analysis}

\subsubsection{Summary Statistics}

For each parameter value, we compute:
\begin{itemize}
    \item Mean, median, standard deviation
    \item 95\% confidence intervals
    \item Interquartile range
\end{itemize}

\subsubsection{Hypothesis Testing}

We employ:
\begin{itemize}
    \item \textbf{Welch's t-test}: For pairwise comparisons (handles unequal variances via Satterthwaite approximation for degrees of freedom)
    \item \textbf{One-way ANOVA}: For comparing means across $\geq 3$ categorical parameters
    \item \textbf{Post-hoc comparisons}: Following significant ANOVA ($p < 0.05$), we perform pairwise Welch's t-tests with BH correction (not Tukey HSD, as BH provides better control of false discovery rate for our many-metric comparisons)
    \item \textbf{Linear regression}: For continuous parameter relationships, reporting $R^2$, slope, and $F$-statistics
    \item \textbf{Kruskal-Wallis $H$}: Non-parametric alternative when normality assumptions are violated (|skewness| $> 1$)
\end{itemize}

Significance level $\alpha = 0.05$. For multiple comparisons, we apply the \textbf{Benjamini-Hochberg} (BH) procedure to control the false discovery rate (FDR) at $q = 0.05$. BH is preferred over Bonferroni correction as it provides greater statistical power while controlling for the expected proportion of false discoveries~\cite{benjamini1995}.

\subsubsection{Effect Size}

We report Cohen's $d$ for pairwise comparisons:
\begin{equation}
d = \frac{\bar{x}_1 - \bar{x}_2}{s_{\text{pooled}}}, \quad s_{\text{pooled}} = \sqrt{\frac{(n_1-1)s_1^2 + (n_2-1)s_2^2}{n_1+n_2-2}}
\end{equation}

For ANOVA, we report both eta-squared ($\eta^2$) and omega-squared ($\omega^2$):
\begin{equation}
\eta^2 = \frac{SS_{\text{between}}}{SS_{\text{total}}}, \quad \omega^2 = \frac{SS_{\text{between}} - df_{\text{between}} \cdot MS_{\text{within}}}{SS_{\text{total}} + MS_{\text{within}}}
\end{equation}

Omega-squared provides a less biased estimate of population effect size than eta-squared, particularly for small samples~\cite{cohen1988statistical}. We interpret effect sizes using Cohen's conventions: $|d| < 0.2$ negligible, $0.2 \leq |d| < 0.5$ small, $0.5 \leq |d| < 0.8$ medium, $|d| \geq 0.8$ large. For $\eta^2$: $< 0.01$ negligible, $0.01$--$0.06$ small, $0.06$--$0.14$ medium, $\geq 0.14$ large.

\subsection{Validation}

\subsubsection{Internal Validity}

We ensure internal validity through:
\begin{itemize}
    \item Deterministic seeding for reproducibility
    \item Multiple runs per configuration (minimum 25; 30 for most experiments)
    \item Systematic parameter variation
\end{itemize}

\subsubsection{External Validity}

We validate simulator outputs against empirical DAO data from DeepDAO, Messari governance reports, and academic literature. Table~\ref{tab:validation} compares baseline simulation results (N=2{,}000 runs) with observed values from real DAOs.

\begin{table}[htbp]
\centering
\caption{Validation: Simulator outputs vs. empirical DAO data}
\label{tab:validation}
\small
\begin{tabular}{@{}lcccc@{}}
\toprule
Metric & Simulator & 95\% CI & Empirical Range & Source \\
\midrule
Proposal Pass Rate & 0.071 & [0.069, 0.072] & [0.30, 0.85] & DeepDAO \\
Average Turnout & 0.059 & [0.058, 0.059] & [0.01, 0.15] & Messari \\
Quorum Reach Rate & 0.227 & [0.225, 0.229] & [0.20, 0.70] & DeepDAO \\
Voter Participation & 0.048 & [0.047, 0.048] & [0.01, 0.10] & \cite{barbereau2022dao} \\
Token Gini & 0.467 & [0.465, 0.470] & [0.70, 0.95] & \cite{nadler2020} \\
Whale Influence & 0.553 & [0.551, 0.554] & [0.40, 0.80] & \cite{fritsch2022} \\
Capture Risk & 0.312 & [0.311, 0.313] & [0.15, 0.55] & \cite{barbereau2022dao} \\
Delegate Conc. & 0.313 & [0.312, 0.314] & [0.30, 0.70] & DeepDAO \\
\bottomrule
\end{tabular}
\end{table}

The simulator achieves 75\% validation (6/8 metrics within empirical ranges). Two notable deviations:
\begin{itemize}
    \item \textbf{Proposal Pass Rate} (7.1\% vs 30-85\%): Our baseline models high quorum failure rates consistent with immature DAOs; mature DAOs achieve higher pass rates through delegate systems.
    \item \textbf{Token Gini} (0.47 vs 0.70-0.95): Our power-law distribution ($\alpha=1.5$) produces moderate inequality; real DAOs often exhibit more extreme concentration.
\end{itemize}

These deviations are addressed in calibration experiments (Section~\ref{sec:results}).

Limitations on external validity:
\begin{itemize}
    \item Agent models are simplified representations
    \item Limited adversarial modeling (attack/defense in experiment~07 only)
    \item Static preferences (no learning)
\end{itemize}

We discuss these limitations in Section~\ref{sec:limitations}.

\subsection{Computational Resources}

% AUTO-UPDATED: computation stats
Experiments were run on:
\begin{itemize}
    \item Hardware: Consumer workstation (details in Appendix)
    \item Concurrency: 4-8 parallel simulations
    \item Total compute time: See reproducibility manifest
    \item Total simulation runs: \totalruns{}
\end{itemize}
